\uuid{UxC1}
\exo7id{757}
\titre{exo7 757}
\auteur{gourio}
\organisation{exo7}
\datecreate{2001-09-01}
\isIndication{false}
\isCorrection{false}
\chapitre{Fonctions circulaires et hyperboliques inverses}
\sousChapitre{Fonctions hyperboliques et hyperboliques inverses}
\cpge{mpsi-mp2i.fonctions-valeurs-reelles-complexes.c, pcsi.techniques-calcul-fonctions.c, ptsi.techniques-calcul-fonctions.c, mpsi-mp2i.nombres-complexes.e, pcsi.nombres-complexes.e, ptsi.nombres-complexes.e, mpsi-mp2i.calcul-algebrique-trigonometrie.d, pcsi.calcul-algebrique-trigonometrie.d, ptsi.calcul-algebrique-trigonometrie.d}
\module{Analyse}
\niveau{L1}
\difficulte{}

\contenu{
\question{
Soit $f:\Rr^{2}\rightarrow \Rr^{2}$ d\'{e}finie par:
$$\forall (x,y)\in \Rr^{2},f(x,y)=(\cos x+\ch y,\cos x\ch y).$$
Discuter et d\'{e}terminer selon $p\in \Rr$ l'image r\'{e}ciproque de $(4,p)$. 
On exprimera $y$ \`{a} l'aide d'un logarithme.
D\'{e}terminer num\'{e}riquement cette image r\'{e}ciproque si $p=-2$.
}
}
